\section{Experiments}
\label{sec:exp}

This paper focuses on the \emph{mechanism} behind Refusal Direction Transplant: whether an alignment direction extracted from an external aligned LLM carries a functional, causal effect on the action-decoding distribution of a vision-language-action model, and how that effect depends on source direction, layer, rank, injection target, and scale. Our evaluation is accordingly built around activation-space and action-token-distribution proxies rather than closed-loop simulator rollouts. We defer rollout-based and adversarial-patch evaluations, as well as a full defense comparison, to a companion study; see Limitations.

\subsection{Setup}
\label{sec:exp:setup}

\paragraph{Backbone.}
We use \textsc{OpenVLA-7b} (\texttt{openvla/openvla-7b}). The external aligned LLM from which the refusal direction is extracted is \texttt{NousResearch/\allowbreak Llama-2-7b-chat-hf}, a non-gated mirror with weights identical to Meta's Llama-2-7b-chat. \textsc{OpenVLA-OFT} serves as a scheme-specific negative control and is discussed in the cross-model analysis (Sec.~\ref{sec:analysis:crossmodel}).

\paragraph{Evaluation data.}
We assemble an \textbf{instruction-conditioned AdvBench split} of 128 harmful instructions drawn from a curated embodied-hazard taxonomy (violence, child safety, fire/electric, chemical, property damage, negligence), paired with visual scenes sampled from \textsc{LIBERO}. We use \textsc{LIBERO} only as a \emph{source of visual contexts} for the VLA prompt, not as a rollout environment: each harmful prompt is paired with a plausible tabletop scene so that the VLA receives the multimodal input distribution it expects. For benign controls we pair Alpaca-drawn benign instructions with the same set of scenes. Unless otherwise stated, ablations are computed over $n{=}32$ harmful prompts with $n_{\text{func}}{=}20$ used for generation-based metrics, matching the \texttt{sanity\_v3c} run.

\paragraph{Attack model.}
This version considers only \emph{textual jailbreak}: harmful instructions delivered in plain language via the \textsc{AdvBench} split above. Image-patch attacks (UADA / UPA / TMA) and white-box adaptive attacks (PGD on $|h_L[\text{last}] \cdot r_L|$) are \emph{out of scope} for this paper and are flagged in Limitations.

\paragraph{Scope disclosure.}
We do not report rollout-based BSR or rollout-derived HRR in this version; the activation-space proxies defined in Sec.~\ref{sec:exp:metrics} suffice for the mechanistic claims of this paper (see Limitations for the resulting deferrals). Concretely, we drop the \textsc{LIBERO-Long} rollout suite, the \textsc{SafeAgentBench} embodied-hazard subset, image-patch attacks, and the PGD adaptive attack from the evaluation scope. What remains is a focused mechanistic characterization of \method{} and \method{}+.

\paragraph{\method{} / \method{}+ configuration.}
Unless an ablation varies these, we use layer $L{=}10$ (selected by the \texttt{last\_text} projection AUC criterion of Sec.~\ref{sec:method:diagnostic}) with $\alpha{=}1.0$. Base \method{} targets action positions only; \method{}+ uses the text+action dual-phase injection of Sec.~\ref{sec:method:dualphase} with the canonical asymmetric ratio $\alpha_{\text{text}}{=}0.3\,\alpha$ and $\alpha_{\text{act}}{=}\alpha$. The rank-$k$ variant (Sec.~\ref{sec:method:rankk}) reports $k \in \{1,3,5,10\}$.

\subsection{Metrics: activation-space refusal proxies}
\label{sec:exp:metrics}

Because we do not execute closed-loop rollouts, we adopt three \emph{simulator-agnostic refusal proxies} computed directly on the VLA's generated action tokens. OpenVLA discretizes each of its 7 action DoFs into 256 bins; bin 128 corresponds to zero motion, and we define a zero-motion window $\mathcal{Z} = [120, 136]$ of $\pm 8$ bins around 128.

\begin{itemize}
  \item \textbf{\texttt{zero\_motion\_mass}} (higher is better): the softmax probability mass the VLA assigns to $\mathcal{Z}$ at each of the 7 action-token positions, averaged over DoFs and prompts. This is a soft, distributional proxy for ``the robot's policy has been pushed toward stopping''.
  \item \textbf{HRR-strict} (higher is better): the fraction of harmful prompts whose \emph{argmax} action tokens all 7 fall in $\mathcal{Z}$. This discretizes the soft proxy into a single-prompt refusal decision.
  \item \textbf{HRR-partial} (higher is better): the fraction of harmful prompts with $\geq 4$ of 7 argmax action tokens in $\mathcal{Z}$. This credits partial refusals (e.g.\ halting the arm but not the gripper) and is more informative when the defense is strong but imperfect.
\end{itemize}

All three are computed via \texttt{generate()}-based decoding of 7 action tokens per prompt, following the \texttt{sanity\_v3c} protocol. We use \texttt{zero\_motion\_mass} as the primary metric in the ablations below because it is continuous and the lowest-variance of the three; the discretized HRR-strict and HRR-partial are defined here for completeness, track \texttt{zero\_motion\_mass} monotonically on the samples we have inspected, and are reported in tabulated form alongside \texttt{zero\_motion\_mass} whenever a headline configuration is discussed.

\paragraph{Roadmap and mapping to contributions.}
The rest of this section establishes the paper's mechanistic claims through a focused ablation grid over the four hyperparameters of Sec.~\ref{sec:method:rdt}--\ref{sec:method:rankk}: (i) direction source vs.\ random control, (ii) injection target, (iii) injection layer, and (iv) subspace rank. Specifically, Sec.~\ref{sec:exp:diagnostic} establishes the diagnostic precondition that the Llama-chat refusal direction carries above-chance linear signal on at least one OpenVLA layer, supporting the structural-safety-gap framing of contribution C1. Sec.~\ref{sec:exp:alpha-base} demonstrates that a single-hook, training-free application of this direction produces a measurable and bounded shift toward zero-motion for both base \method{} and the dual-phase \method{}+ (C2; \method{}+ at $\alpha{=}1$ is the headline $+2.84$\,pp configuration). Sec.~\ref{sec:exp:rand} contains the paper's causal centerpiece: the gain is attributable to the specific transplanted geometry rather than to generic hidden-state perturbation (C3). Sec.~\ref{sec:exp:target} isolates the contribution of position-aware targeting (C4, and the specific advantage of \method{}+ over base \method{} and over blanket \textsc{VLM-Guard}-style injection). Sec.~\ref{sec:exp:rankk} interrogates whether action-space refusal is single- or multi-dimensional, yielding preliminary C5 evidence under the $\sqrt{k}$-normalized reading. Sec.~\ref{sec:exp:layer} characterizes the layer-sensitivity of the intervention, and also surfaces a diagnostic limitation---per-layer AUC and Cohen's $d$ are not the sole determinants of intervention efficacy---which we treat as future work rather than as a new contribution.

\subsection{Diagnostic: per-layer AUC and Cohen's $d$}
\label{sec:exp:diagnostic}

Before intervening, we ask where in \textsc{OpenVLA}'s residual stream the Llama-chat refusal direction carries discriminative signal. Following Sec.~\ref{sec:method:diagnostic}, we compute, at each candidate layer $L$, the projection $p = h_L \cdot r_L$ and report the binary-classification AUC of $p$ for harmful vs.\ benign inputs at both the last text token and the first action token, together with the Cohen's $d$ effect size. Full per-layer curves are plotted in Fig.~\ref{fig:mechanism}(a); the headline numbers are as follows. Layer 10 is the natural operating point with \texttt{last\_text} AUC $=0.638$ and \texttt{first\_action} AUC $=0.589$ (Cohen's $d$ $=+0.36$ and $+0.27$), both positive and close---consistent with the structural-safety-gap claim of Sec.~\ref{sec:intro} that harm/benign information is present in the residual stream but the action-decoding path does not act on it. Layer 18 is the maximally \emph{anti}-aligned deep layer: Cohen's $d$ $=-1.10$ at \texttt{last\_text}, while a 5-fold out-of-fold linear probe still attains AUC $>0.95$ at both positions across all layers tested. The harm/benign information is thus retained under \textsc{OpenVLA}'s 27-epoch behavior-cloning fine-tuning, but the one-dimensional Llama-chat axis has been progressively reorganized along the action-decoding path. We select $L{=}10$ for subsequent intervention ablations.

\subsection{\method{} and \method{}+: the $\alpha$ sweep}
\label{sec:exp:alpha-base}

We inject the rank-1 refusal direction at $L{=}10$ and sweep $\alpha \in \{0, 0.5, 1, 2, 5\}$ for both variants: base \method{} (action-only target) and \method{}+ (text+action dual-phase with $\alpha_{\text{text}}{=}0.3\,\alpha,\alpha_{\text{act}}{=}\alpha$ from Sec.~\ref{sec:method:dualphase}). The baseline (no intervention) attains $\texttt{zero\_motion\_mass} = 0.2274$ and HRR-partial $= 0.050$.

\begin{table}[t]
\centering\small
\begin{tabular}{r r r r r}
\toprule
 & \multicolumn{2}{c}{base \method{} (action)} & \multicolumn{2}{c}{\method{}+ (text+action)} \\
\cmidrule(lr){2-3}\cmidrule(lr){4-5}
$\alpha$ & \texttt{z\_mass} & $\Delta$ & \texttt{z\_mass} & HRR-p \\
\midrule
0.0 & 0.2274 & --- & 0.2274 & 0.050 \\
0.5 & 0.2351 & $+0.008$ & 0.2446 & 0.100 \\
\textbf{1.0} & 0.2441 & $+0.017$ & \textbf{0.2558} & \textbf{0.150} \\
2.0 & \textbf{0.2468} & $\mathbf{+0.019}$ & 0.2479 & 0.100 \\
5.0 & 0.2415 & $+0.014$ & 0.2297 & 0.000 \\
\bottomrule
\end{tabular}
\caption{$\alpha$ sweep for base \method{} and \method{}+ at $L{=}10$ (\texttt{sanity\_v3c}, \texttt{sanity\_v4}; $n_{\text{func}}{=}20$). \method{}+ at $\alpha{=}1$ is the headline configuration: $+2.84$\,pp \texttt{zero\_motion\_mass} over baseline and HRR-partial tripled from $0.050$ to $0.150$. Both variants show a monotone-then-regressing profile (plateau at $\alpha \in [1,2]$, collapse at $\alpha{=}5$)---evidence of out-of-distribution extrapolation when hidden states are pushed too far along a direction whose semantics were fit on Llama-chat's training distribution. HRR-strict is $0$ throughout at $n{=}20$: a $+2.84$\,pp soft shift is not enough to push all 7 DoFs into $\mathcal{Z}$ simultaneously.}
\label{tab:alpha-plus}
\end{table}

At $\alpha{=}1$, \method{}+ beats base \method{} by $+1.2$\,pp \texttt{zero\_mass} and $+5$\,pp HRR-partial---the generation-phase-aware dual injection of Sec.~\ref{sec:method:dualphase} adds real value on top of action-only targeting. The HRR-strict flat-zero row is informative in itself: a single-hook, norm-bounded intervention shifts the action distribution but does not ``turn off'' all 7 DoFs at once, which is consistent with the training-free positioning of \method{}+.

\subsection{Direction-source control ($\Delta_{\text{rand}}$)}
\label{sec:exp:rand}

To test whether the gain above is attributable to the \emph{transplanted safety geometry} rather than to generic hidden-state perturbation, we replace $r_L$ with a random unit vector of the same norm (fixed seed) and repeat the intervention. Define $\Delta_{\text{rand}} = \texttt{zero\_mass}(r_L) - \texttt{zero\_mass}(\text{rand})$; positive $\Delta_{\text{rand}}$ is the core evidence for contribution C3. At $L{=}10$, action-only, we obtain $\Delta_{\text{rand}} = +0.0157$ at $\alpha{=}1$ (refusal $0.2441$ vs.\ random $0.2284$) and $+0.0146$ at $\alpha{=}2$ (refusal $0.2468$ vs.\ random $0.2322$); at $\alpha{=}5$ the refusal direction underperforms random ($0.2415$ vs.\ $0.2620$, $\Delta_{\text{rand}}{=}-0.0205$), consistent with out-of-distribution extrapolation. We accordingly restrict the effective operating range to $\alpha \in [0.5, 2]$; the full curves appear in Fig.~\ref{fig:mechanism}(b). This is the narrowest causal claim the experiment can support and is the paper's mechanistic centerpiece: within a bounded $\alpha$ range, \emph{which} direction is transplanted matters, and the Llama-chat refusal axis carries content that a random axis of the same norm does not.

\subsection{Target-position ablation}
\label{sec:exp:target}

At $L{=}10, \alpha{=}1$ with $r_L$ held fixed, we vary the injection target and obtain \texttt{zero\_mass}: \textbf{text+action} ($\textbf{0.2558}$, \method{}+) $>$ \textbf{action} ($0.2441$, base \method{}) $>$ \textbf{all} ($0.2410$, \textsc{VLM-Guard}-style blanket injection) $>$ \textbf{text} ($0.2372$). Two takeaways. First, \method{}+ beats base \method{} by $+1.2$\,pp \texttt{zero\_mass}: the phase-aware dual injection of Sec.~\ref{sec:method:dualphase} adds real value on top of action-only targeting. Second, the \emph{all} condition---where every position is perturbed---underperforms both \method{} variants, undercutting the naive intuition that more intervention is better and giving a clean differentiation from \textsc{VLM-Guard}-style untargeted injections.

\subsection{Rank-$k$ subspace sweep}
\label{sec:exp:rankk}

We next ask whether the harm/benign distinction lives in a single direction or a low-rank subspace. Following Sec.~\ref{sec:method:rankk}, we replace $r_L$ with $\sum_{j=1}^{k} v_j$ where $v_j$ are the top right singular vectors of the contrast matrix, and sweep $k \in \{1,3,5,10\}$ at $L{=}10, \alpha{=}1$, action-only targeting.

\begin{table}[t]
\centering\small
\begin{tabular}{r r r}
\toprule
$k$ & unnorm.\ \texttt{z\_mass} & $\sqrt{k}$-norm.\ \texttt{z\_mass} \\
\midrule
1  & 0.2249 & 0.2249 \\
3  & 0.2110 & 0.2328 \\
\textbf{5}  & 0.2209 & \textbf{0.2358} \\
10 & 0.2455 & 0.2291 \\
\bottomrule
\end{tabular}
\caption{Rank-$k$ subspace sweep, $L{=}10, \alpha{=}1$, action-only. Unnormalized $r_k = \sum_{j \le k} v_j$ (\texttt{sanity\_v3c}) injects $\sqrt{k}\times$ more norm as $k$ grows, confounding any gain with an effective $\alpha$ increase; $\sqrt{k}$-normalized $r_k = (1/\sqrt{k})\sum_{j\le k} v_j$ (\texttt{sanity\_v4}) holds norm constant. Under the norm-controlled column, $k{=}5$ peaks at $0.2358$ (beating rank-1 by $+1.1$\,pp) and the unnormalized $k{=}10$ advantage disappears.}
\label{tab:rankk}
\end{table}

The $\sqrt{k}$-normalized column of Table~\ref{tab:rankk} is the correct lens for contribution C5: rank-3 and rank-5 improve over rank-1 by $+0.8$ and $+1.1$\,pp \texttt{zero\_mass} respectively, with rank-5 peak and a collapse by $k{=}10$ (consistent with the top-10 SVD directions reaching past the meaningful harm/benign subspace into near-null components). We accordingly read C5 as a \emph{small but positive} multi-dimensional gain: action-space refusal is not exactly rank-1, but it is effectively low-rank (3--5).

\subsection{Layer sensitivity}
\label{sec:exp:layer}

Sweeping $L \in \{8,10,12,14,16,18,20\}$ at $\alpha{=}1$, action-only (sanity\_v3c), we find \texttt{zero\_motion\_mass} is nearly flat across $L \in \{10, 14, 16, 18\}$ (all within $\pm 0.3$\,pp of $L{=}10$'s $0.2441$), despite the diagnostic AUC and $d$ varying dramatically over the same range (Fig.~\ref{fig:mechanism}a). Per-layer linear-probe signal is thus not the sole determinant of intervention efficacy: \texttt{zero\_motion\_mass} is primarily sensitive to the magnitude of perturbation injected, and OpenVLA's action-BC training reorganizes the deep representations such that a one-dimensional probe is not the right lens on it. We treat this as a diagnostic limitation, pick $L{=}10$ for the convenience of positive alignment at both positions, and leave a principled layer-selection rule to future work.

\paragraph{Summary.}
Across Tables~\ref{tab:alpha-plus}--\ref{tab:rankk} the consistent picture is: a transplanted refusal direction produces a small but non-trivial, saturating shift of OpenVLA's action distribution toward zero-motion; the shift is specific to the refusal axis within a bounded $\alpha$ range (Sec.~\ref{sec:exp:rand}); it is maximized when injected at both text and action positions (Sec.~\ref{sec:exp:target} and Table~\ref{tab:alpha-plus}, $+2.84$\,pp); and it is insensitive to the exact intervention layer within $L \in [10, 18]$. Rollout-based BSR/HRR, composition with other defenses, and adaptive-attack robustness are deferred to the companion study.
